\begin{figure}[ht]
\centering
\begin{tikzpicture}[
  >=Stealth,
  dipolo/.style={->, thick, brown!70!black},
  Pvec/.style={->, thin, purple!70!black},
  leader/.style={dashed, gray, thin},
  every node/.style={font=\small}
]

% --- Regione esterna: continuo polarizzato ---
\fill[purple!15, rounded corners=4pt] (-4.2,-3.2) rectangle (4.2,3.2);

% Frecce di P sparse (continuo uniformemente polarizzato)
\draw[Pvec] (-3.6,-2.6) -- ++(0,0.5);
\draw[Pvec] (-3.6,-0.5) -- ++(0,0.5);
\draw[Pvec] (-3.6, 1.8) -- ++(0,0.5);
\draw[Pvec] (-2.2, 2.4) -- ++(0,0.5);
\draw[Pvec] ( 0.0, 2.6) -- ++(0,0.5);
\draw[Pvec] ( 2.2, 2.4) -- ++(0,0.5);
\draw[Pvec] ( 3.6, 1.8) -- ++(0,0.5);
\draw[Pvec] ( 3.6,-0.5) -- ++(0,0.5);
\draw[Pvec] ( 3.6,-2.6) -- ++(0,0.5);
\draw[Pvec] ( 2.2,-2.8) -- ++(0,0.5);
\draw[Pvec] ( 0.0,-2.8) -- ++(0,0.5);
\draw[Pvec] (-2.2,-2.8) -- ++(0,0.5);

% --- Sfera di Lorentz ---
\draw (0,0) circle (1.6);
\draw[dashed, thick] (0,0) circle (1.6);

% Dipoli discreti dentro la sfera (orientazioni varie)
\draw[dipolo] (-0.9, 0.6) -- ++(0.05,0.45);
\draw[dipolo] ( 0.8, 0.7) -- ++(0.10,0.45);
\draw[dipolo] (-0.8,-0.5) -- ++(-0.10,0.45);
\draw[dipolo] ( 0.9,-0.4) -- ++(0.10,0.45);
\draw[dipolo] ( 0.0, 1.0) -- ++(0.05,0.40);
\draw[dipolo] (-1.1, 0.0) -- ++(0.00,0.45);
\draw[dipolo] ( 1.1, 0.0) -- ++(0.00,0.45);
\draw[dipolo] (-0.3,-1.0) -- ++(-0.05,0.45);
\draw[dipolo] ( 0.4,-1.0) -- ++(0.05,0.45);

% --- Punto centrale: la molecola di interesse ---
\fill[red!80!black] (0,0) circle (1.8pt);

% --- Etichette con linee guida ---
\draw[leader] (1.35, 1.0) -- (4.2, 2.6);
\node[anchor=west] at (4.2, 2.6) {Continuo, $ \v{E}_P $};

\draw[leader] (0.7, 0.5) -- (4.2, 1.0);
\node[anchor=west] at (4.2, 1.0) {Dipoli discreti, $\mathbf{E}_{\text{micro}}$};

\draw[leader] (0, 0) -- (4.2,-0.6);
\node[anchor=west] at (4.2,-0.6) {Punto di interesse, $ \v{E}_\textnormal{loc} $};

\end{tikzpicture}
\caption{Sfera mesoscopica con cui modelliamo gli effetti di campo locale: tutto il mezzo è modellizzato tramite polarizzazione dando 
il campo $ \v{E}_P $, mentre dentro abbiamo il contributo dei dipoli che contiamo come $ \v{E}_{\text{micro}} $.}
\label{figure:_effetti_campo_locale}
\end{figure}